% Demo slide showcasing definition/example/caveat highlight blocks. Copy
% from here when starting a new slide; see README.md for the
% naming/versioning convention.
%
% `definitionblock` added (2026-07-06) alongside the original three, same
% day it was introduced in Slidemaster_preamble_uni-graz.tex for real content on
% slides/4uq-Definition-of-Capital-Market-Efficiency-current.tex and
% slides/dny-Leasing-Definition-current.tex -- every reusable block/frame
% type in this deck gets a demo here, so this one shouldn't be the
% exception.
%
% v1-1-0 (2026-09-18, direct instruction): `infoblock` demo added -- same
% "every reusable block type gets a demo here" reasoning as definitionblock
% above; this one had been in real use (e.g.
% slides/k4p-Introductory-Question-Answer-current.tex's own
% `gothamfloatinginfoblock`, which just wraps `infoblock` in a minipage for
% corner-overlay placement) without ever getting its own entry here. Shown
% as a plain in-flow block, matching how the other four in this demo are
% shown -- the floating/corner-overlay positioning is `gothamfloatinginfoblock`'s
% own concern, not `infoblock`'s own visual style, which is all this slide
% demos. Five stacked full-width blocks no longer fit one frame (confirmed
% by a first compile: `Overfull \vbox 60.06pt too high`, the fifth block
% silently pushed off the bottom of the visible page with no on-page trace
% -- not merely cramped) -- split into two `columns` instead of stacking
% all five full-width.
%% Included in decks: test-slides
%%

\setSlideVersion{1-1-0}
\begin{standardframe}{Highlighting a definition or formula}
	\footnotesize
	\begin{columns}[t,onlytextwidth]
		\begin{column}{0.48\linewidth}
			\begin{definitionblock}{Definition, distinguished by $\triangleq$}
				Use \texttt{definitionblock} for a definition that should stand
				out from an ordinary \texttt{block} at a glance.
			\end{definitionblock}
			\begin{block}{Definition}
				Use \texttt{block} for neutral definitions.
			\end{block}
			\begin{exampleblock}{Example}
				Use \texttt{exampleblock} for worked examples.
			\end{exampleblock}
		\end{column}
		\begin{column}{0.48\linewidth}
			\begin{alertblock}{Caveat}
				Use \texttt{alertblock} to flag caveats or warnings.
			\end{alertblock}
			\begin{infoblock}{Info}
				Use \texttt{infoblock} (or \texttt{gothamfloatinginfoblock} to
				place it in a corner, see \texttt{k4p}) for supplementary
				information.
			\end{infoblock}
		\end{column}
	\end{columns}
\end{standardframe}
